% Canonical Paper IV section.
Paper~I defines a marked spinal history as a chronological word of one-hole
arithmetic contexts.  We recall only the structure needed to discuss quotient
fibers and continuation.  Throughout this section $K$ is a field, and all
projective contexts are assumed non-degenerate.

\subsection{Free histories and two semantic categories}

Let $\Sigma_K$ be a selected set of invertible one-hole arithmetic contexts.
It may contain translations, nonzero scalings, sign changes, and inversion.
A free history is a word
\[
  \gamma=s_1s_2\cdots s_n,
  \qquad s_i\in\Sigma_K,
\]
listed in chronological order.  Its induced transformation is
\begin{equation}
  \rho(\gamma)=\rho(s_n)\cdots\rho(s_2)\rho(s_1)
  \in G:=\PGL_2(K).
  \label{eq:chronological-evaluation}
\end{equation}
Thus a newly appended future context acts by left multiplication.  This
chronology convention is essential in Section~\ref{sec:condensation-continuation}.

\begin{definition}[Ordinary and projective evaluation]
For an initial value $x\in K$, ordinary evaluation $\nu_x(\gamma)$ is defined
only when every intermediate arithmetic operation is defined in $K$.
Projective evaluation is the everywhere-defined action of $\rho(\gamma)$ on
$\PP^1(K)$.  We write
\[
  \operatorname{dom}_{\mathrm{ord}}(\gamma)\subseteq K
\]
for the ordinary domain.
\end{definition}

For example, the projective continuation $z\mapsto c/z$ sends $0$ to
$\infty$, while ordinary division by zero is undefined.  The projective target
therefore does not erase the ordinary-domain label.  Later residual definitions
will assign a distinguished observation $\bot$ to an illegal ordinary
continuation.

\begin{definition}[Word-labelled history category]
Let $\cH_K$ have objects $q\in\PP^1(K)$.  A morphism from $q$ to $q'$ is a
pair $(\gamma,q)$ consisting of a literal context word with
$\rho(\gamma)q=q'$.  If $\gamma$ is followed chronologically by $\delta$,
the composite is labelled by the concatenated history $\gamma\star\delta$
and satisfies
\[
  \rho(\gamma\star\delta)=\rho(\delta)\rho(\gamma).
\]
Formal inverses of context symbols give a groupoid completion
$\cH_K^{\mathrm{inv}}$.  Projective evaluation defines a functor
\[
  \rho:\cH_K^{\mathrm{inv}}
  \longrightarrow
  \PP^1(K)\sslash G,
\]
to the action groupoid by forgetting the literal word and retaining its
projective operator.  Distinct word-labelled arrows are not identified merely
because they have the same image in $G$.
\end{definition}

This construction is intentionally modest.  Associativity, commutativity,
distributivity, and program transformations may later be added as rewrite
$2$-cells between word-labelled arrows.  They are not imposed as literal
history equality.

\subsection{The equality tower}

A bounded history $(x,\gamma)$ supports several non-equivalent equality tests.
For compatible histories $\gamma$ and $\delta$, consider:
\begin{align*}
  \gamma&=\delta &&\text{literal history equality},\\
  \rho(\gamma)&=\rho(\delta) &&\text{projective operator equality},\\
  \rho(\gamma)q&=\rho(\delta)q &&\text{equality at a projective state $q$},\\
  \nu_x(\gamma)&=\nu_x(\delta) &&\text{ordinary endpoint equality},\\
  \chi(\gamma)&=\chi(\delta) &&\text{equality of a selected charge shadow}.
\end{align*}
Only the forward implications forced by the selected maps are automatic.
Endpoint equality at one input does not imply operator equality; operator
equality does not imply history equality; and charge equality need not imply
operator equality.  These failures are not pathologies.  They are the fibers
that the present paper studies.

\begin{proposition}[Nested kernel pairs]
Let $X\xrightarrow{f}Y\xrightarrow{g}Z$ be maps.  The kernel pair of $f$ is
contained in the kernel pair of $g\circ f$:
\[
  f(x)=f(x')\quad\Longrightarrow\quad g(f(x))=g(f(x')).
\]
The converse holds exactly when $g$ is injective on $f(X)$.
\end{proposition}
\begin{proof}
The forward implication is functoriality of equality.  If $g$ is injective on
$f(X)$, equality after $g$ implies equality before $g$.  Conversely, if the
converse implication holds for all $x,x'$, then $g$ is injective on the image
of $f$.
\end{proof}

This proposition is the set-theoretic skeleton of the full tower.  Its
computational interpretation requires a future class, developed in
Section~\ref{sec:contextual-residuals}.

\subsection{Group-valued relative defects}

Endpoint subtraction is natural in an affine chart but not projectively
invariant.  The operator group supplies a comparison that survives inversion.

\begin{definition}[Relative projective defect]
For invertible histories $\gamma$ and $\delta$, define
\begin{equation}
  \cK(\gamma,\delta)
  \defeq \rho(\delta)^{-1}\rho(\gamma)\in G.
  \label{eq:relative-projective-defect}
\end{equation}
\end{definition}
It satisfies
\[
  \cK(\gamma,\gamma)=1,
  \qquad
  \cK(\gamma,\delta)^{-1}=\cK(\delta,\gamma),
\]
and the cocycle identity
\begin{equation}
  \cK(\gamma,\eta)
  =\cK(\delta,\eta)\cK(\gamma,\delta).
  \label{eq:defect-cocycle}
\end{equation}

\begin{proposition}[Affine torsion as a coordinate of the relative defect]
Suppose
\[
  \rho(\gamma)(x)=sx+A_\gamma,
  \qquad
  \rho(\delta)(x)=sx+A_\delta,
  \qquad s\ne0.
\]
Then
\[
  \cK(\gamma,\delta)(x)=x+\frac{A_\gamma-A_\delta}{s}.
\]
The ordinary endpoint defect $A_\gamma-A_\delta$ is therefore the translation
coordinate of the projective relative defect, multiplied by the common scale.
\end{proposition}
\begin{proof}
The inverse of $x\mapsto sx+A_\delta$ is
$x\mapsto s^{-1}(x-A_\delta)$.  Composition with the first affine map gives
the formula.
\end{proof}

Once the two operators do not have the same affine linear part, a scalar
endpoint difference depends on the initial value.  Once inversion is admitted,
it may also cross poles.  The group-valued defect remains defined at the
operator level, while its endpoint observation still requires domain data.

\subsection{Relations, neutral histories, and fibers}

\begin{definition}[Operator relation and neutrality]
An ordered pair $(\gamma,\delta)$ is an operator relation when
$\rho(\gamma)=\rho(\delta)$.  A history $\gamma$ is operator-neutral when
$\rho(\gamma)=1$.  It is charge-neutral relative to a charge map $\chi$ when
$\chi(\gamma)$ equals the charge of the empty history.
\end{definition}

Operator-neutral histories form the kernel of $\rho$ only when the history
language has been given a compatible group structure.  In the free monoid they
form a submonoid preimage of the identity; in the formal groupoid they form
isotropy arrows above the relevant object.  Charge-neutral and
operator-neutral histories need not coincide.

\begin{example}[Four elementary separations]
The following examples use affine maps and are proved by direct composition.
\begin{enumerate}
  \item $x\mapsto x+1$ and $x\mapsto2x$ agree at $x=1$ but are different
  operators.
  \item The words ``add $1$, then add $2$'' and ``add $2$, then add $1$'' are
  distinct histories inducing the same map $x\mapsto x+3$.
  \item ``add $p$, then scale by $e^q$'' and the opposite order have the same
  ACS endpoint but differ by translation $p(e^q-1)$.
  \item ``add $p$, then add $-p$'' is a nonempty operator-neutral history.
\end{enumerate}
\end{example}

The operator fiber $\rho^{-1}(g)$ and the endpoint fiber
$\{\gamma:\nu_x(\gamma)=y\}$ are therefore distinct objects.  The first is
independent of $x$; the second need not be.  Paper~IV uses ``history fiber''
only after naming the map whose fiber is intended.

\subsection{Chronological end versus operand slot}

The side on which a future operator is multiplied is not the operand-slot
chirality of an arithmetic context.  Operand-slot chirality records whether
the evolving value occupies the first or second input of one binary operation.
Chronological side records whether a new operator is composed before or after
an existing product.  In the convention \eqref{eq:chronological-evaluation}, a
future context acts on the left regardless of its operand slot.  A theorem
relating slot mirror, matrix transpose, inverse, or right continuation would
require additional structure and is not assumed here.
